\subsubsection{二元函数的极值}

\begin{define}[Local Maximum]
    Suppose $f$ is defined on an open neighborhood $D\subseteq\mathbb{R}^n$, and let $x^0\in D$. If there exists $B(x^0,\delta)\subseteq D$ such that $f(x)\leq f(x^0)$ for all $x\in B(x^0,\delta)$, then $x^0$ is called a local maximal point of $f$, and $f(x^0)$ is called a local maximum.
\end{define}

\begin{define}[Critical Point]
    If $f$ is differentiable at $x^0$ and $df|_{x^0}=0$, then $x^0$ is called a \emph{critical point} (驻点) of $f$.
\end{define}

Fact: If $x^0$ is a local maximal point of $f$ and $f$ is differentiable at $x^0$, then $x^0$ is critical, i.e. $df|_{x^0}=0$ or $\nabla f|_{x^0}=0$.

\begin{pf}
    For any $v\in\mathbb{R}^n$, $t=0$ is a local maximal point of $f(x^0+tv)\triangleq\varphi(t)$. Since $\varphi$ is differentiable at $t=0$, we have $\varphi'(0)=0$, i.e. $df|_{x^0}(v)=\frac{\partial f}{\partial v}|_{x^0}=\langle\nabla f|_{x^0}, v\rangle=0$. Since $v$ is arbitrary, $df|_{x^0}=0$.
\end{pf}

\begin{example}
    $f(x,y)=x^2+y^2$, $g(x,y)=x^2-y^2$, $(0,0)$ is critical, but $(0,0)$ is local minimum for $f$ but not for $g$
\end{example}

$$\left.\begin{pmatrix}
        f_{xx} & f_{xy} \\
        f_{yx} & f_{yy}
    \end{pmatrix}\right|_{(0,0)}=\begin{pmatrix}
        2 & 0 \\
        0 & 2
    \end{pmatrix}\quad\left.\begin{pmatrix}
        g_{xx} & g_{xy} \\
        g_{yx} & g_{yy}
    \end{pmatrix}\right|_{(0,0)}=\begin{pmatrix}
        2 & 0  \\
        0 & -2
    \end{pmatrix}
$$

Suppose $f$ is twice differentiable at $0\in\mathbb{R}^n$ and $0$ is a critical point. Then
$$
    f(x)=f(0)+\frac{1}{2}Q(x)+R(x),
$$
where $Q(x)=\sum_{i,j=1}^n\frac{\partial^2f}{\partial x_i\partial x_j}(0)x_ix_j$, and $\left|\frac{R(x)}{|x|^2}\right|\to 0$ as $x\to 0$.

The matrix $\operatorname{Hess}(f)|_0=\left(\frac{\partial^2f}{\partial x_i\partial x_j}(0)\right)_{1\leq i,j\leq n}$ is called the Hessian matrix of $f$ at $0$.

\begin{thm}
    Suppose $\frac{\partial f}{\partial x_i}$ is differentiable at $x^0$ for $1\leq i\leq n$, and $\left.\frac{\partial f}{\partial x_i}\right|_{x^0}=0$ for $i=1,2,\dots,n$.
    \begin{enumerate}
        \item If $\operatorname{Hess}(f)|_{x^0}$ is positive (negative) definite, then $x^0$ is a local minimal (maximal) point of $f$;
        \item If $\operatorname{Hess}(f)|_{x^0}$ has rank $n$ (sometimes called non-degenerate 非退化), but is not definite, i.e. it has both positive and negative eigenvalues, then $x^0$ is a \emph{saddle point} 鞍点
    \end{enumerate}
\end{thm}

\begin{pf}
    $f(x)=f(x^0)+\frac{1}{2}(x-x_0)^T\operatorname{Hess}(f)|_{x^0}(x-x_0)+R(x-x^0)$, where $\frac{R(x-x^0)}{|x-x^0|^2}\to 0$ as $x\to x^0$.

    $$
        \implies\frac{f(x)-f(x^0)}{|x-x^0|^2}=\frac{1}{2}\frac{(x-x_0)^T\operatorname{Hess}(f)|_{x^0}(x-x^0)}{|x-x_0|^2}+\frac{R(x-x^0)}{|x-x^0|^2}.
    $$

    Assume $x^0=0$ and consider the behavior of $\frac{x^THx}{|x|^2}$. Since it is homogeneous of degree $0$,
\end{pf}

